\documentclass[12pt,a4paper]{amsart}

\usepackage{amsmath,amssymb,amsthm}
\usepackage{geometry}
\usepackage{hyperref}
\usepackage{booktabs}
\usepackage{array}
\usepackage{microtype}
\usepackage{enumitem}
\usepackage{xcolor}

\geometry{margin=2.4cm, top=2.8cm, bottom=3cm}

\hypersetup{
  colorlinks=true,
  linkcolor=blue!70!black,
  citecolor=blue!70!black,
  urlcolor=blue!70!black
}

%% ---- Theorem Environments ----
\theoremstyle{plain}
\newtheorem{theorem}{Theorem}[section]
\newtheorem{lemma}[theorem]{Lemma}
\newtheorem{corollary}[theorem]{Corollary}
\newtheorem{proposition}[theorem]{Proposition}
\theoremstyle{definition}
\newtheorem{definition}[theorem]{Definition}
\newtheorem{postulate}[theorem]{Postulate}
\newtheorem{conjecture}[theorem]{Conjecture}
\theoremstyle{remark}
\newtheorem{remark}[theorem]{Remark}

%% ---- Shortcuts ----
\newcommand{\vp}{\varphi}
\newcommand{\Sthree}{S^3}
\newcommand{\lP}{\ell_{\mathrm{P}}}
\newcommand{\Hzero}{H_0}
\newcommand{\Heff}{H_{\mathrm{eff}}}
\newcommand{\trec}{t_{\mathrm{rec}}}
\newcommand{\Rrec}{R_{\mathrm{rec}}}
\newcommand{\Lam}{\ensuremath{\Lambda}}
\newcommand{\LCDM}{\ensuremath{\Lambda}CDM}
\newcommand{\half}{\tfrac{1}{2}}
\newcommand{\Rfrac}{R_{\mathrm{frac}}}
\newcommand{\Rlight}{R_{\mathrm{light}}}
\newcommand{\Vtot}{V_{\mathrm{total}}}
\newcommand{\Vobs}{V_{\mathrm{obs}}}
\newcommand{\chiL}{\chi_{\mathrm{L}}}
\newcommand{\rhotot}{\rho_{\mathrm{total}}}
\newcommand{\rhoobs}{\rho_{\mathrm{obs}}}
\newcommand{\Eeff}{E_{\mathrm{eff}}}

%% ================================================================
\title[Geometric--Causal--Spectral Cosmology on Compact $S^3$]{
  Geometric--Causal--Spectral Cosmology\\
  on a Compact Three-Sphere:\\
  $S^3$ Kernel, Causal Light Front, Density Tracking,\\
  and Redshift as a Density-Gradient Integral}

\author{Andrei Preece}
\address{Independent Researcher, Berlin, Germany}
\email{andre.preece@web.de}

\author{Boris Batenin}
\address{Independent Researcher, Berlin, Germany}
\email{batenin@web.de}

\date{April 2026}

\subjclass[2020]{83F05, 83F30, 11B39, 58J50, 83C45}
\keywords{compact $S^3$ topology, FLRW $k=+1$, $S^3$ kernel,
  causal light front, observable volume, two-sector expansion,
  Fibonacci hierarchy, golden ratio, Perron--Frobenius eigenvalue,
  density tracking, redshift gradient, observable Hubble law,
  Hubble tension}

\begin{document}

\begin{abstract}
\textbf{Positioning.} This paper proposes a geometrically constrained
phenomenological framework on compact $\Sthree$. It is \emph{not} a
replacement of $\Lambda$CDM, \emph{not} a first-principles derivation
of cosmological parameters, and \emph{not} a modification of local
Einstein gravity. The central conceptual content is: the observable
universe should be distinguished from the full spatial geometric
volume, and density may track this observable domain dynamically rather
than fill it instantaneously.

We formulate a Geometric--Causal--Spectral cosmology on a compact,
simply-connected three-sphere $\Sthree(R(t))$ as the joint evolution
of three fundamental fields: the geometric radius $R(t)$; the
dimensionless causal light angle
$\chiL(t)=\int_0^t c\,dt'/R(t')$; and the energy density $\rho(t)$.

We introduce the $\Sthree$ kernel
$K(\chi)=2\sin^2\!\chi/(\chi-\sin\chi\cos\chi)$ governing radial volume
growth, with limits $K\sim 3/\chi$ (Euclidean) and $K(\pi)=0$
(antipodal). The observable volume is
$\Vobs=2\pi R^3[\chiL-\half\sin 2\chiL]$, distinct from
$\Vtot=2\pi^2 R^3$.

The radius obeys a sigmoid-bridged two-sector law with a discrete
Fibonacci pre-recombination phase (Perron--Frobenius eigenvalue $\vp$)
and an exponential post-recombination phase governed by an
\emph{effective} expansion rate $\Heff\neq\Hzero$.

The observable density obeys the postulated relaxation
$\dot\rhoobs=-\Gamma(t)[\rhoobs-\rhotot]$ with $\Gamma=c/(R\chiL)$.
The cosmological redshift is the line integral
$\ln(1+z)=\int_0^\ell p(V(r))K(\chi(r))/R(r)\,dr$.
The observable Hubble parameter is window-dependent,
$H(\ell)=3cp/\ell$, suggesting a phenomenological two-regime structure with ratio
$\Hzero^{\mathrm{late}}/\Hzero^{\mathrm{CMB}}\approx 1.08$--$1.10$,
consistent with the observed CMB--late-time split at the
phenomenological level (not a parameter-free prediction).

We do not eliminate $\Lam$, do not derive $\Hzero$, $\Lam$, or the
vacuum energy density from first principles, do not replace the
Friedmann equations, and do not provide a theory of quantum
gravity. The Hubble-tension consistency is a phenomenological
check, not a derivation.
\end{abstract}

\maketitle
\tableofcontents

%% =================================================================
\section{Geometry of $S^3$}
\label{sec:geom}
%% =================================================================

We assume spatial sections $\Sigma$ that are compact, connected,
orientable, boundaryless, with trivial fundamental group, and
endowed with a round Riemannian metric. By the Poincar\'e--Perelman
theorem~\cite{Perelman2002,Perelman2003a,Perelman2003b},
$\Sigma\cong\Sthree$ uniquely.

The total geometric volume of $\Sthree(R)$ is
\begin{equation}
  \Vtot(R) = 2\pi^2 R^3.
  \label{eq:Vtot}
\end{equation}
The geodesic radial coordinate $\chi\in[0,\pi]$ measures dimensionless
angular distance from a chosen pole. The volume of a geodesic ball
of comoving angle $\chi$ on $\Sthree(R)$ is
\begin{equation}
  \Vobs(\chi, R) = 2\pi R^3\bigl[\chi - \half\sin(2\chi)\bigr].
  \label{eq:Vobs}
\end{equation}
At $\chi=\pi$, $\Vobs=\Vtot$; for $\chi<\pi$, $\Vobs<\Vtot$.

\medskip
\begin{center}
\fbox{\quad$\Vobs \neq \Vtot$\quad}
\end{center}

%% =================================================================
\section{The $S^3$ Kernel}
\label{sec:kernel}
%% =================================================================

\begin{definition}[$\Sthree$ kernel]
  \label{def:K}
  The $\Sthree$ kernel governing radial volume growth is
  \begin{equation}
    K(\chi) \equiv \frac{1}{V}\frac{dV}{d\chi}
       = \frac{2\sin^2\!\chi}{\chi - \sin\chi\cos\chi},
    \quad \chi\in(0,\pi).
    \label{eq:Kdef}
  \end{equation}
\end{definition}

\begin{theorem}[Limits of $K$]
  \label{thm:Klimits}
  $K(\chi)\to 3/\chi$ as $\chi\to 0$;
  $K(\chi)\to 0$ as $\chi\to\pi$;
  $K(\pi/2) = 4/\pi\approx 1.273$.
\end{theorem}

\begin{proof}
  Taylor expansion: $\sin\chi\cos\chi=\sin(2\chi)/2=\chi-2\chi^3/3+\cdots$,
  so denominator $\chi-\sin\chi\cos\chi\to 2\chi^3/3$. Numerator
  $2\sin^2\!\chi\to 2\chi^2$. Ratio $\to 3/\chi$.
  At $\chi=\pi$: $\sin\pi=0$ in numerator, $\pi$ in denominator.
  At $\chi=\pi/2$: numerator $2$, denominator $\pi/2$.
\end{proof}

\begin{remark}[Physical interpretation]
  $K(\chi)$ is the fractional rate of volume growth per unit angular
  distance. In small-$\chi$ Euclidean regime, $K\sim 3/\chi$ recovers
  the standard $4\pi r^2/[(4\pi/3)r^3]=3/r$ shell-to-volume ratio.
  Near the antipode, $K\to 0$: no further volume is gained as the
  geodesic ball wraps around the closed manifold.
\end{remark}

%% =================================================================
\section{Three Fundamental Fields}
\label{sec:fields}
%% =================================================================

\begin{definition}[Three fundamental fields]
  \label{def:fields}
  The cosmological state is specified by three independent fields:
  \begin{enumerate}[label=\textbf{(F\arabic*)},noitemsep]
  \item \emph{Geometric radius} $R(t)>0$: curvature radius of $\Sthree$
    at cosmic time $t$.
  \item \emph{Causal light angle}
    $\chiL(t) = \int_0^t c\,dt'/R(t')\in[0,\pi]$, the dimensionless
    angular distance traversed by a null geodesic.
  \item \emph{Energy density} $\rho(t)>0$: distributed on $\Sthree(R(t))$.
  \end{enumerate}
\end{definition}

\begin{remark}[Why $\chiL$ rather than a dimensionful $r_L$]
  $\chiL$ is the natural causal coordinate on $\Sthree$: it is
  dimensionless, bounded above by $\pi$, and parametrises the causal
  geometry without conflating spatial growth with light propagation.
  The proper causal distance is $r_L=R(t)\chiL(t)$.
\end{remark}

%% =================================================================
\section{Observable Volume}
\label{sec:obsvol}
%% =================================================================

\begin{definition}[Observable volume]
  \label{def:Vobs_t}
  The observable volume at time $t$ is the geodesic ball of angle
  $\chiL(t)$:
  \begin{equation}
    \Vobs(t) = 2\pi R^3(t)\bigl[\chiL(t) - \half\sin 2\chiL(t)\bigr].
    \label{eq:Vobs_t}
  \end{equation}
\end{definition}

\begin{theorem}[Properties of $\Vobs(t)$]
  \label{thm:Vobsprop}
  For $\chiL\ll 1$: $\Vobs\approx (4\pi/3)(R\chiL)^3$ (Euclidean limit).
  For $\chiL=\pi$: $\Vobs=\Vtot$.
  Today, $\chiL\sim 0.07$ and $\Vobs/\Vtot\sim 7\times 10^{-5}$.
\end{theorem}

\begin{proof}
  Direct: small-$\chi$ Taylor of \eqref{eq:Vobs_t} gives $(4\pi/3)R^3\chi^3$.
  Numerical evaluation at $\chi=0.07$, $R=1.85\times 10^{27}\,\mathrm{m}$
  gives $\Vobs/\Vtot=7.4\times 10^{-5}$.
\end{proof}

%% =================================================================
\section{Fractal Spectral Structure}
\label{sec:phi}
%% =================================================================

\begin{lemma}[Minimality of the Fibonacci recursion]
  \label{lem:minimal}
  Among all positive integer second-order recursions
  $a_{n+1}=p\,a_n+q\,a_{n-1}$ ($p,q\in\mathbb{Z}_{\geq 1}$), the
  Fibonacci recursion ($p=q=1$) has the smallest Perron--Frobenius
  growth rate $\lambda_+=\vp$.
\end{lemma}

\begin{proof}
  Companion matrix
  $T_{p,q}=\bigl(\begin{smallmatrix}p&q\\1&0\end{smallmatrix}\bigr)$
  has dominant eigenvalue $\lambda_+(p,q)=(p+\sqrt{p^2+4q})/2$,
  minimised at $p=q=1$ with $\lambda_+=\vp$.
\end{proof}

\begin{theorem}[$\varphi$ as Perron--Frobenius eigenvalue]
  \label{thm:phi}
  The Fibonacci ratio $F_{n+1}/F_n$ converges to $\vp=(1+\sqrt5)/2$,
  with rate $|F_{n+1}/F_n-\vp|=O(\vp^{-2n})$.
\end{theorem}

\begin{proof}
  Standard via Binet's formula~\cite{Hardy1979,Meyer2000}.
\end{proof}

\begin{remark}[Status of $\varphi$]
  The golden ratio is the dominant eigenvalue of the minimal positive
  second-order recursion (Lemma~\ref{lem:minimal}). The $\vp$-ladder
  supplies the discrete spectral hierarchy used to organise scales;
  its observational realisation in the cosmological dynamics remains
  model-dependent.
\end{remark}

\begin{corollary}[Spectral bridge]
  \label{cor:bridge}
  If the discrete pre-recombination radii satisfy $F_n/R_n=\mathrm{const}$,
  then $R_n\propto\vp^n$ and density scales as $\rho_n\propto\vp^{-3n}$.
\end{corollary}

%% =================================================================
\section{Two-Sector Radius Law (Centre of the Framework)}
\label{sec:twosector}
%% =================================================================

\begin{definition}[Two-sector radius law]
  \label{def:twosector}
  Let $\trec>0$ be the recombination time, $\Delta t>0$ its width,
  and $\sigma(t)=[1+e^{-(t-\trec)/\Delta t}]^{-1}$ the sigmoid weight.
  Define
  \begin{align}
    \Rfrac(t)  &= \lP\,\vp^{N(t)},
      &&\text{(fractal sector)}\label{eq:Rfrac}\\
    \Rlight(t) &= \Rrec\,e^{\Heff(t-\trec)},
      &&\text{(light sector)}\label{eq:Rlight}
  \end{align}
  with $\Rrec=\Rfrac(\trec)=\Rlight(\trec)$. The cosmological radius is
  \begin{equation}
    \boxed{
      R(t) = \Rfrac(t)\bigl[1-\sigma(t)\bigr] + \Rlight(t)\,\sigma(t).
    }
    \label{eq:twosector}
  \end{equation}
\end{definition}

\begin{theorem}[$C^1$-smoothness]
  \label{thm:C1}
  Under matching $\Rfrac(\trec)=\Rlight(\trec)$: $R(t)$ and $\dot R(t)$
  are continuous everywhere; $R(\trec)=\Rrec$;
  $\dot R(\trec)=\half[\dot{\Rfrac}+\dot{\Rlight}]_{\trec}$.
\end{theorem}

\subsection{Effective Expansion Rate $\Heff\neq H_0$}

\begin{definition}[Effective post-recombination rate]
  \label{def:Heff}
  \begin{equation}
    \Heff = \frac{\ln(R_{\mathrm{today}}/\Rrec)}{t_0-\trec}.
    \label{eq:Heff}
  \end{equation}
\end{definition}

\begin{remark}[$\Heff\neq H_0$]
  $\Heff$ is the \emph{average} expansion rate over the post-recombination
  era; $H_0$ is the present-day \emph{instantaneous} rate.
  With $R_{\mathrm{today}}/\Rrec\approx 1100$ and $t_0-\trec\approx 4.35\times 10^{17}\,\mathrm{s}$:
  \begin{equation}
    \Heff\approx 1.61\times 10^{-17}\,\mathrm{s}^{-1}\approx 7.4\,H_0,
    \label{eq:Heffvalue}
  \end{equation}
  with $H_0=2.18\times 10^{-18}\,\mathrm{s}^{-1}$ (Planck~2018).
  These are different physical quantities. $\Heff$ is a
  matching parameter of the two-sector radius law, not the measured
  late-time Hubble constant. Earlier
  presentations~\cite{Preece2025FBS3R_GoldenRatio} that used
  $H=2.45\times 10^{-17}\,\mathrm{s}^{-1}$ in the exponential were
  approximating $\Heff$, but mis-identifying it as $H_0$.
\end{remark}

%% =================================================================
\section{Causal Light Propagation}
\label{sec:lightprop}
%% =================================================================

\begin{theorem}[Null-geodesic equation on $\Sthree$]
  \label{thm:null}
  In the FLRW metric on $\Sthree$, null geodesics satisfy
  $c\,dt = R(t)\,d\chi$, integrating to
  \begin{equation}
    \chiL(t) = \int_0^t \frac{c\,dt'}{R(t')}.
    \label{eq:chiLdef}
  \end{equation}
\end{theorem}

\subsection{Light as the Generator of Observable Volume}

The dimensionless causal angle $\chiL(t)$ defines, via~\eqref{eq:Vobs_t},
the observable volume. Light is therefore not merely an in-sphere
propagating signal: \emph{it is the operative agent that determines
what is causally accessible} at time $t$.

In the pre-recombination era ($\sigma\ll 1$), the cosmic plasma is
opaque and the operational meaning of $\Vobs$ is restricted. In the
post-recombination era, free-streaming photons make $\Vobs(t)$ the
operationally meaningful domain.



%% =================================================================
\section{Compatibility with FLRW and Einstein Gravity}
\label{sec:flrw}
%% =================================================================

The framework is compatible with general relativity. Spatial topology
$\Sthree$ corresponds to FLRW geometry with curvature index $k=+1$:
\begin{equation}
  ds^2 = -c^2dt^2 + a^2(t)R_0^2\,d\Omega_3^2.
  \label{eq:flrw_metric}
\end{equation}
Local Einstein equations are unchanged; the Friedmann constraint reads
\begin{equation}
  H_{\mathrm{GR}}^2 = \frac{8\pi G}{3}\rho - \frac{c^2}{a^2 R_0^2}
    + \frac{\Lam c^2}{3}.
  \label{eq:friedmann}
\end{equation}

\begin{remark}[Compatibility, not replacement]
  The present framework is compatible with this geometry but does
  \emph{not} claim to replace the Friedmann equations. The Gradient
  Redshift Ansatz (Postulate~\ref{post:gradient}) is a coarse-grained
  observational ansatz, not a substitute for local geodesic optics
  in GR. The framework introduces no local modification of Einstein
  gravity.
\end{remark}

%% =================================================================
\section{Density Dynamics: A Postulate of Causal Tracking}
\label{sec:density}
%% =================================================================

\begin{definition}[Two density fields]
  \label{def:rhotwo}
  Assuming an effective conserved energy $\Eeff$:
  \begin{align}
    \rhotot(t) &= \frac{\Eeff}{\Vtot(R(t))}\quad\text{(global)}\label{eq:rhotot}\\
    \rhoobs(t) &= \frac{E_{\mathrm{eff,within}\,\Vobs}(t)}{\Vobs(t)}
                 \quad\text{(observable)}\label{eq:rhoobs}
  \end{align}
\end{definition}

\begin{postulate}[Causal Density Tracking]
  \label{post:tracking}
  The observable density relaxes towards the global density at a rate
  set by the inverse causal information time:
  \begin{equation}
    \boxed{
      \dot\rhoobs(t) = -\Gamma(t)\bigl[\rhoobs(t) - \rhotot(t)\bigr],
      \qquad \Gamma(t) = \frac{c}{R(t)\,\chiL(t)}.
    }
    \label{eq:tracking}
  \end{equation}
\end{postulate}

\begin{remark}[Status of Postulate~\ref{post:tracking}]
  This equation models relaxation of observable density toward the
  global geometric density. It is \emph{not} derived from microscopic
  kinetic theory in this paper. The form $\Gamma=c/(R\chiL)$ is the
  simplest causal scaling consistent with the framework. A microscopic
  derivation is deferred to future work. Physically, the relaxation
  rate $\Gamma=c/(R\chiL)$ can be interpreted as a finite propagation
  speed of equilibration across the causal domain, bounded by light
  travel.
\end{remark}

\begin{remark}[Density tracks volume, not the reverse]
  The postulate formalises the central thesis: density does not
  instantaneously fill the geometric volume; it tracks the growth of
  the causal observable region with finite relaxation time
  $\Gamma^{-1}=R\chiL/c$.
\end{remark}

%% =================================================================
\section{Redshift as a Density-Gradient Integral: Phenomenological Ansatz}
\label{sec:redshift}
%% =================================================================

\subsection{Phenomenological Gradient Redshift Law}

\begin{postulate}[Gradient Redshift Ansatz]
  \label{post:gradient}
  The photon angular frequency $\omega$ along a null geodesic
  parametrised by affine parameter $\lambda$ evolves according to a
  local volume gradient:
  \begin{equation}
    \frac{d\ln\omega}{d\lambda} = -p(V)\,\frac{1}{V}\frac{dV}{d\lambda},
    \label{eq:omega_local}
  \end{equation}
  where $p(V)$ is a dimensionless coarse-grained morphology index
  characterising effective photon coupling to local volume change.
\end{postulate}

\begin{remark}[Status of Postulate~\ref{post:gradient}]
  Equation~\eqref{eq:omega_local} is \emph{not} derived from local
  GR or QFT. It is introduced as an effective coarse-grained transport
  rule for photon frequency along null geodesics in a compact $\Sthree$
  causal domain. This ansatz should be interpreted as an effective
  coarse-grained description and is not unique; alternative transport
  laws may exist.
\end{remark}

\begin{remark}[Status of $p(V)$]
  $p(V)$ is \emph{not} a fluid equation-of-state parameter and must
  not be identified with $1+w$. It is a coarse-grained
  morphology/rarefaction index. We constrain $0<p(V)\leq p_{\max}$
  and adopt a piecewise-constant simplification $p_{\mathrm{rec}},
  p_{\mathrm{late}}$ for distinct observational windows. A
  first-principles derivation of $p(V)$ is deferred to future work.
\end{remark}

\subsection{Redshift--Distance Relation on $\Sthree$}

\begin{proposition}[Conditional consequence of the Gradient Redshift Ansatz]
  \label{thm:redshift}
  \emph{Conditional on Postulate~\ref{post:gradient}}, the redshift
  along a null geodesic from a source at proper distance $\ell$ obeys
  \begin{equation}
    \frac{d\ln(1+z)}{dr} = \frac{p(V)\,K(\chi(r))}{R(r)},
    \label{eq:dlnz}
  \end{equation}
  giving
  \begin{equation}
    \boxed{
      \ln(1+z) = \int_0^\ell \frac{p(V(r))\,K(\chi(r))}{R(r)}\,dr.
    }
    \label{eq:zintegral}
  \end{equation}
\end{proposition}

\begin{proof}
  Combining~\eqref{eq:omega_local} with $1+z=\omega_{\mathrm{em}}/\omega_{\mathrm{obs}}$
  gives $d\ln(1+z)=p\,d\ln V$. Using $d\ln V/dr=K(\chi)/R$ along the
  null geodesic gives~\eqref{eq:dlnz}.
\end{proof}

\begin{remark}[Physical interpretation]
  Redshift is the integral of the local fractional volume gradient
  along the photon's path, weighted by $p$. The $\Sthree$ kernel
  $K(\chi)$ encodes the geometric correction: in the Euclidean limit
  $K\sim 3/\chi$ recovers the standard inverse-distance dilution; near
  the antipode $K\to 0$ photons stop acquiring further redshift.
\end{remark}

%% =================================================================
\section{Observable Hubble Law}
\label{sec:hubble}
%% =================================================================

\begin{proposition}[Effective windowed Hubble slope]
  \label{thm:Hwindow}
  \emph{Conditional on Postulate~\ref{post:gradient}}, for an
  observation window of proper extent $\ell$ with $\ell>\ell_{\min}$,
  approximately constant $p$, and $\chi\ll 1$ (so $K\sim 3R/r$),
  the redshift--distance relation can be approximated as
  $\ln(1+z)\approx 3p\ln(\ell/\ell_0)$. For $z\ll 1$, the
  effective windowed Hubble slope can be written as
  \begin{equation}
    \boxed{
      H(\ell) = \frac{cz}{\ell}\Big|_{z\ll 1} = \frac{3cp}{\ell}.
    }
    \label{eq:Hell}
  \end{equation}
  The $\ell\to 0$ limit is not physical.
\end{proposition}

\begin{proof}
  In the small-$\chi$ regime $K(\chi)/R=3/r$, so~\eqref{eq:zintegral}
  gives $\ln(1+z)=3p\int_{\ell_0}^\ell dr/r = 3p\ln(\ell/\ell_0)$.
  Linearising for $z\ll 1$ gives~\eqref{eq:Hell}.
\end{proof}

\begin{remark}[$H(\ell)$ is not a fundamental Hubble parameter]
  $H(\ell)$ is \emph{not} a universal global expansion rate. It is an
  observationally windowed slope obtained after smoothing the local
  $\Sthree$ gradient over a finite scale $\ell$. The expression is
  valid only for finite observational windows $\ell>\ell_{\min}$;
  the $\ell\to 0$ limit is not physical. The dependence on $\ell$
  implies that observational inference of $\Hzero$ may be sensitive
  to window selection effects.
\end{remark}

%% =================================================================
\section{Two Regimes and the Hubble Tension}
\label{sec:tworegimes}
%% =================================================================

The window-dependent Hubble law naturally accommodates two distinct effective values:

\begin{align}
  \Hzero^{\mathrm{CMB}} &= \frac{3c\,p_{\mathrm{rec}}}{\ell_{\mathrm{rec}}},
    \quad p_{\mathrm{rec}}\approx 1/3,\label{eq:HCMB}\\
  \Hzero^{\mathrm{late}} &= \frac{3c\,p_{\mathrm{loc}}}{\ell_{\mathrm{loc}}}.
    \label{eq:Hlate}
\end{align}

\begin{theorem}[Window ratio]
  \label{thm:Hratio}
  \begin{equation}
    \frac{\Hzero^{\mathrm{late}}}{\Hzero^{\mathrm{CMB}}}
      = \frac{p_{\mathrm{loc}}}{p_{\mathrm{rec}}}\cdot
        \frac{\ell_{\mathrm{rec}}}{\ell_{\mathrm{loc}}}.
    \label{eq:Hratio}
  \end{equation}
\end{theorem}

For the parameters consistent with the two-sector law, this ratio
falls in the range
\begin{equation}
  \frac{\Hzero^{\mathrm{late}}}{\Hzero^{\mathrm{CMB}}}\approx 1.08\text{--}1.10,
\end{equation}
consistent at the percent-level with the observed CMB--late-time split between Planck
($\sim 67.4\,\mathrm{km/s/Mpc}$) and SH0ES local
measurements~\cite{Riess2022} ($\sim 73\,\mathrm{km/s/Mpc}$):
ratio $\approx 73.0/67.4=1.083$. This consistency is a phenomenological check, not a parameter-free prediction.

\begin{remark}[Status of the Hubble-tension consistency]
  The ratio $\Hzero^{\mathrm{late}}/\Hzero^{\mathrm{CMB}}\approx 1.08$--$1.10$
  should be interpreted as a phenomenological consistency check, not
  as a parameter-free prediction. $p$ and $\ell$ are taken from
  observational fits, not derived from first principles. We do not
  claim that the Hubble tension is solved. This mechanism does
  not exclude standard systematic explanations of the Hubble tension.
\end{remark}

%% =================================================================
\section{Numerical Consistency}
\label{sec:numerics}
%% =================================================================

\begin{table}[h]
\centering
\caption{Numerical consistency check. Planck~2018 parameters
  used throughout~\cite{Planck2018}.}
\label{tab:numerics}
\begin{tabular}{lcc}
\toprule
Quantity & Symbol & Value \\
\midrule
Planck length & $\lP$ & $1.616\times 10^{-35}\,\mathrm{m}$ \\
Recombination time & $\trec$ & $1.20\times 10^{13}\,\mathrm{s}$ \\
Recombination radius & $\Rrec$ & $9.46\times 10^{20}\,\mathrm{m}$ \\
Today (cosmic) & $t_0$ & $4.35\times 10^{17}\,\mathrm{s}$ \\
Today scale factor & $R(t_0)$ & $\sim 1.04\times 10^{24}\,\mathrm{m}$ \\
Effective rate & $\Heff$ & $1.61\times 10^{-17}\,\mathrm{s}^{-1}$ \\
Present-day rate & $H_0$ & $2.18\times 10^{-18}\,\mathrm{s}^{-1}$ \\
Ratio $\Heff/H_0$ & --- & $\approx 7.4$ \\
Curvature radius & $R_{\mathrm{curv}}$ & $\geq 1.85\times 10^{27}\,\mathrm{m}$ \\
Today $\chiL$ & --- & $\sim 0.07$ \\
$\Vobs/\Vtot$ today & --- & $\sim 7\times 10^{-5}$ \\
$N(\trec)$ Fibonacci level & --- & $\approx 266.85$ \\
$K(\pi/2)$ kernel mid-value & --- & $4/\pi\approx 1.273$ \\
$\Hzero^{\mathrm{late}}/\Hzero^{\mathrm{CMB}}$ predicted & --- & $1.08$--$1.10$ \\
$\Hzero^{\mathrm{late}}/\Hzero^{\mathrm{CMB}}$ observed & --- & $1.083$ \\
\bottomrule
\end{tabular}
\end{table}

%% =================================================================
\section{Physical Interpretation}
\label{sec:interp}
%% =================================================================

The framework rests on three operational principles:

\begin{enumerate}[label=(\arabic*),noitemsep]
\item \textbf{Light defines what is accessible.} The observable volume
  $\Vobs(t)$ is the geodesic ball of comoving angle $\chiL(t)$,
  bounded above by $\pi$ on compact $\Sthree$.
\item \textbf{Density is dynamical, not instantaneous.} $\rhoobs$
  relaxes toward $\rhotot$ at rate $\Gamma=c/(R\chiL)$
  (Postulate~\ref{post:tracking}).
\item \textbf{Redshift is a path integral.} $\ln(1+z)$ is the line
  integral of $p(V)K(\chi)/R$ along null geodesics
  (Theorem~\ref{thm:redshift}); $H(\ell)=3cp/\ell$ is therefore
  window-dependent.
\end{enumerate}

The Hubble tension between CMB and local measurements is a direct
consequence: different observational ladders correspond to different
effective $(p,\ell)$ pairs, giving distinct effective Hubble parameters
without contradiction.

%% =================================================================
\section{Scope and Limitations}
\label{sec:scope}
%% =================================================================

\subsection{What This Paper Does}

\begin{enumerate}[label=(Y\arabic*),noitemsep]
\item Establishes the three-field framework: $R(t)$, $\chiL(t)$, $\rho(t)$.
\item Introduces the $\Sthree$ kernel $K(\chi)$ with its limits.
\item Derives $\vp$ as the Perron--Frobenius eigenvalue
  (Theorem~\ref{thm:phi}).
\item States the $C^1$-smooth two-sector law with $\Heff\neq H_0$.
\item Postulates causal density tracking
  (Postulate~\ref{post:tracking}).
\item Reformulates redshift conditionally as a density-gradient integral
  (Theorem~\ref{thm:redshift}).
\item Predicts a window-dependent Hubble law $H(\ell)=3cp/\ell$.
\item Identifies $\Hzero^{\mathrm{late}}/\Hzero^{\mathrm{CMB}}\approx 1.08$--$1.10$,
  consistent with the observed Hubble tension.
\end{enumerate}

\subsection{What This Paper Does Not Do}

\begin{enumerate}[label=(N\arabic*),noitemsep]
\item Does \emph{not} derive $\Lam$ from first principles.
\item Does \emph{not} derive $\Hzero$ from first principles.
\item Does \emph{not} derive $p(V)$ from microscopic theory.
\item Does \emph{not} replace the Friedmann equations.
\item Does \emph{not} provide a full quantum gravity model.
\item Does \emph{not} claim that the Hubble tension is solved.
\item Does \emph{not} claim that the observed vacuum density is
  predicted without calibration.
\item Does \emph{not} eliminate dark energy. The exponential sector
  requires $\Lam$ as in \LCDM.
\item Does \emph{not} derive $\Heff$, $\Eeff$, or
  $\Gamma=c/(R\chiL)$ from first principles. They are matching
  parameters and postulates.
\item Does \emph{not} provide a wave function of the universe;
  earlier ansatze~\cite{Preece2025FBS3R_GoldenRatio} are withdrawn.
\end{enumerate}

%% =================================================================
\section*{Conclusions}
%% =================================================================

The principal contributions are:
\begin{enumerate}[noitemsep]
\item $\Sthree$ kernel: $K(\chi)\to 3/\chi$ (Euclidean), $K(\pi)=0$
  (antipodal).
\item Volume distinction: $\Vobs/\Vtot\sim 7\times 10^{-5}$ today.
\item $\vp$ as Perron--Frobenius eigenvalue, not a parameter.
\item Two-sector $C^1$-smooth law with $\Heff\approx 7.4\,H_0$.
\item Causal density tracking
  $\dot\rhoobs=-[c/(R\chiL)][\rhoobs-\rhotot]$ as postulate.
\item Conditional reformulation of redshift as gradient integral
  $\ln(1+z)=\int p(V)K(\chi)/R\,dr$.
\item Window-dependent Hubble law $H(\ell)=3cp/\ell$.
\item Estimated within the present phenomenological framework: Hubble tension ratio $\approx 1.08$--$1.10$, consistent with
  observation $1.083$.
\end{enumerate}

\medskip
\noindent\textbf{Final positioning.} The main result of this paper
is \emph{not} a replacement of \LCDM, \emph{not} a derivation of
cosmological parameters, and \emph{not} a modification of local GR.
It is a controlled geometric reinterpretation of observable
cosmological volume and density on compact $\Sthree$. The distinction
between $\Vobs$ and $\Vtot$, together with the causal density-tracking
postulate and the gradient redshift ansatz, provides a phenomenological
framework for studying whether observed expansion parameters may depend
on finite observational windows.

\medskip
\begin{center}
\textit{All rigorous content lies in the geometry: $\Sthree$ topology,
volume formula, kernel $K(\chi)$, Perron--Frobenius~$\vp$.\\
All physical interpretation is at the level of postulates, ans\"atze,
and phenomenological consistency.}
\end{center}

%% =================================================================
\begin{thebibliography}{99}

\bibitem{Preece2025FBS3R_GoldenRatio}
A.~Preece and B.~Batenin,
Fractal Bindu $(S^3)$ of Reality: The Golden Ratio as the Principle of
Perfection and Harmony in Interdependent Existence,
Zenodo (2025), \href{https://doi.org/10.5281/zenodo.15538573}{doi:10.5281/zenodo.15538573}.

\bibitem{Planck2018}
Planck Collaboration,
Planck 2018 results.\ VI.\ Cosmological parameters,
Astron.\ Astrophys.\ \textbf{641}, A6 (2020).

\bibitem{Perelman2002}
G.~Perelman,
The Entropy Formula for the Ricci Flow,
arXiv:math/0211159 (2002).

\bibitem{Perelman2003a}
G.~Perelman,
Ricci Flow with Surgery on Three-Manifolds,
arXiv:math/0303109 (2003).

\bibitem{Perelman2003b}
G.~Perelman,
Finite Extinction Time for the Solutions to the Ricci Flow,
arXiv:math/0307245 (2003).

\bibitem{Hardy1979}
G.~H.~Hardy and E.~M.~Wright,
\textit{An Introduction to the Theory of Numbers}, 5th ed.,
Oxford Univ.\ Press (1979).

\bibitem{Meyer2000}
C.~D.~Meyer,
\textit{Matrix Analysis and Applied Linear Algebra}, SIAM (2000).

\bibitem{Riess2022}
A.~G.~Riess et~al.,
A Comprehensive Measurement of the Local Value of the Hubble Constant,
Astrophys.\ J.\ Lett.\ \textbf{934}, L7 (2022).

\bibitem{PreeceBatenin2026Bell}
A.~Preece and B.~Batenin,
Compact Simply-Connected Cosmology and the Geometric Origin of Quantum
Nonlocality (Version 3.0), Berlin (2026).

\end{thebibliography}

\end{document}
